%% Section body, shared verbatim by ../paper/main.tex (preprint) and
%% ../submission/body.tex (Information Sciences submission).  Information Sciences requires a
%% Conclusions section written to be intelligible outside the specialism; this is it.
On a dependent stream, a streaming quasi-Newton method can get the estimate right and the uncertainty
wrong. The error is not small and does not go away: the reported variance omits the autocovariances of the
score, so a nominal $95\%$ interval converges to a coverage below $95\%$ and stays there.
\Cref{prop:coverage} says exactly where --- at $2\Phi(z_{\alpha/2}/\sqrt{\kappa})-1$, computable from one
variance-inflation factor $\kappa$. At the inflation we measure on hourly air quality that is
$\RealPredAirq$, so more than half of the intervals a practitioner would report on that stream miss the
value they are meant to cover.

The remedy is a change of reading order rather than an extra estimation stage. Consume the stream in
contiguous blocks and take one preconditioned step per block. The block-averaged gradients are then, up to
a factor, the non-overlapping batch means of the score process, and batch means estimate exactly the
long-run variance the interval needs. The object the statistics requires and the object the computation
already produces are the same, so the resulting interval is \emph{cheaper} than the invalid one: an
instantaneous plug-in covariance costs $O(d^2)$ per observation, the blocked long-run covariance $O(d^2)$
per block. Gapping the rank-one curvature updates is the second change, and it is a cost knob rather than
an inference device.

Three practical consequences. Block length is the one parameter that needs thought, and needs less than
one might fear: because the blocks carry gradients rather than parameter values, their length is set by
how quickly the \emph{data} forget, not by how quickly the optimiser settles, so $b$ of order $\log N$
suffices where iterate-based methods need a power of $N$. Blocking has a price we pay and report: it makes
a step-size condition non-trivial, and without the $O(d)$ safeguard of \eqref{eq:step} our own method
diverges on a persistent regime-switching design. And honest intervals are wider --- on our designs
$\WidthRatioLo$--$\WidthRatioHi$ times the dependence-blind width, which is $\sqrt{\kappa}$ and not a
tuning choice. The narrower interval was never correct; it was the wrong length.

What we have not done bounds the claims, and \Cref{sec:limits} states each boundary precisely. The
construction itself is general: any streaming method that forms block gradients and maintains a curvature
estimate can report a long-run variance at no extra order of cost.
